\documentclass[11pt]{article}

\usepackage[margin=1in]{geometry}
\usepackage{amsmath,amssymb,amsthm}
\usepackage{booktabs}
\usepackage{tabularx}
\usepackage{array}
\usepackage{enumitem}
\usepackage{microtype}
\usepackage{xcolor}
\usepackage[hidelinks]{hyperref}

\newcolumntype{Y}{>{\raggedright\arraybackslash}X}

\newtheorem{proposition}{Proposition}
\newtheorem{definition}{Definition}
\newtheorem{observation}{Observation}
\newtheorem{counterexample}{Counterexample}

\newcommand{\project}{\mathsf{project}}
\newcommand{\fork}{\mathsf{fork}}
\newcommand{\settle}{\mathsf{settle}}
\newcommand{\effects}{\mathsf{effects}}
\newcommand{\prefix}{\mathsf{prefix}}
\newcommand{\suffix}{\mathsf{suffix}}
\newcommand{\code}[1]{\texttt{#1}}
\newcommand{\artifact}{Agent Runtime Laws}

\title{Confluence and Fork-Safety in Event-Sourced Agent Runtimes}
\author{Yohei Nakajima}
\date{Research draft --- August 2026}

\begin{document}

\maketitle

\begin{abstract}
Event-sourced agent runtimes commonly claim deterministic replay, cheap
forking, and end-to-end lineage.  These claims are attractive but incomplete.
A pure reducer is deterministic only after the event sequence is fixed; a
reactive runtime may admit multiple firing orders; an oracle result must be
captured before it can be replayed; and an event-log fork cannot undo an email,
charge, publication, or model-call cost that already occurred.  This paper asks
when replay and fork claims are sound, rather than proposing another minimal
agent primitive.

We present an executable F\# specification that separates replayable facts from
requested effects, exposes scheduler order as a parameter, and assesses fork
safety relative to four distinct properties: retained-prefix projection,
strict execution replay, continuation in a prefix-reconstructed environment,
and equivalence to a counterfactual external world.  Effects are described by
three orthogonal dimensions: external footprint, replay source, and lifecycle.
FsCheck laws and preserved counterexamples show that settlement is neither
terminating nor confluent in general.  Overlapping writers yield different
normal forms, while a second counterexample shows that disjoint declared writes
remain insufficient when behavior reads and emitted triggers interfere.

We validate the conditions against two runtime artifacts comprising 5,564 run
logs and 314,348 source events.  In a public Synthetic Players capsule, all
317,296 projection cuts are sound under the stated projection observation, but
36,251 cuts split an outstanding model request and are unsound for external
continuation.  The store contains 121 observed forks; their 16,625 shared-prefix
events have zero structural mismatches modulo child run identity and global
database sequence, and all observed cuts avoid external requests.  A separate
24-run offline bridge study has 66.7\% decision agreement but 0\% ordered-path
agreement, illustrating that replay, path, projection, and decision equivalence
are different assertions.  The results turn ``cheap fork'' from a performance
slogan into a property-indexed semantic claim with checkable preconditions.
\end{abstract}

\noindent\textbf{Artifact status.}
The code, law suite, conformance vectors, and evidence ledger accompany this
draft.  The public capsule is revision- and hash-pinned.  The bridge artifact is
locally hash-bound but its containing source directory has no commit or remote;
we therefore treat it as non-reproducible outside the local environment until
published.

\section{Introduction}

Agent runtimes increasingly retain an append-only record of model calls, tool
results, human decisions, state updates, and coordination events.  Event
sourcing makes several desirable properties seem immediate.  If state is a
fold over the log, a run can be replayed.  If a child log shares a prefix with a
parent, a run can be forked without recomputing that prefix.  If causes are
recorded, artifacts can be traced back to model and tool activity.  ActiveGraph
states these properties as consequences of making the log the source of truth
\cite{nakajima2026log}.

This paper formalizes and empirically checks the replay and fork properties
claimed on architectural grounds in Nakajima (2026).  The relationship is
deliberate: the earlier work describes a deployed architecture, while the
present work asks which side conditions make its architectural claims sound.
Negative results about the originating runtime are therefore part of the
contribution rather than exceptions to be hidden.

Two gaps motivate the work.  First, ``deterministic replay'' often means
deterministic replay \emph{after oracle outputs have been fixed}.  This is
useful, but the phrase suppresses operational questions.  Was the model output
captured faithfully?  Is a request retained without its result?  Does replay
serve the recorded result or make a new call?  Is the runtime itself
schedule-independent, or does a particular FIFO order merely produce one
repeatable execution?

Second, event sourcing in a business application normally reconstructs
application state; it does not claim to reverse the external world.  That
distinction becomes load-bearing for agents.  A retrospectively discarded log
suffix may contain an email, a charge, a delete, a publication, or a billable
model call.  The forked trace can omit the effect, but the world cannot be
assumed to be one in which the effect never happened.  A correct fork contract
must say which property it preserves.

The initial version of this project asked whether

\begin{equation}
  A : S \times E \rightarrow S \times F^*
\end{equation}

was the irreducible agent primitive.  That framing is not novel.  The shape is
a Mealy-style transition system \cite{mealy1955method}, closely resembles the
Elm update architecture \cite{elmarchitecture}, and appears in functional
event sourcing through the Decider pattern \cite{chassaing2021decider}.  We
concede the minimality lane and instead exploit the effects/events seam that
becomes visible when agents interact with irreversible external systems.

\paragraph{Contributions.}
This draft makes four contributions.

\begin{enumerate}[leftmargin=*]
  \item It gives an executable operational model in which scheduler order is
  explicit.  The model separates termination, blocking on an external effect,
  and bounded divergence.

  \item It preserves counterexamples showing that general settlement is neither
  terminating nor confluent.  It further shows that pairwise-disjoint
  \emph{declared writes} are insufficient when behavior reads and emitted
  triggers interfere.  A restricted isolated-writer family does converge, but
  we do not claim a minimal general condition.

  \item It defines property-relative fork assessment over three orthogonal
  effect dimensions: footprint, replay source, and lifecycle.  This separates
  exact retained-prefix projection from execution replay, continuation, and a
  counterfactual-world claim.

  \item It checks the laws against 5,564 runtime logs, 314,348 source events,
  every possible cut in each log, and 121 observed forks.  The validation finds
  both a strong positive result---zero structural mismatches in the observed
  shared prefixes---and sharp limitations on which hypothetical cuts license
  continuation or counterfactual claims.
\end{enumerate}

\paragraph{Genre.}
This is not a mechanized-proof paper.  Laws are stated as propositions, checked
over generated families with FsCheck, and evaluated exhaustively over the cuts
of the available logs.  A failed law produces a minimized, named regression
fixture.  Passing checks are finite evidence over the specified generators and
artifacts, not universal proofs.  This choice is intentional: calculi with
machine-checked composition, termination, and information-flow results already
exist \cite{liu2026lambdaa,garby2026llmbda}.  Our comparative advantage is an
executable runtime contract connected to real runtime artifacts.

\paragraph{Roadmap.}
Section~\ref{sec:model} defines the executable model and observations.
Section~\ref{sec:confluence} reports termination and confluence results.
Section~\ref{sec:fork} defines fork properties and effect conditions.
Section~\ref{sec:grades} connects those conditions to replayability grades.
Sections~\ref{sec:method} and~\ref{sec:validation} describe the artifact and
empirical validation.  We then position the work, discuss implications, and
state limitations.

\section{Executable Runtime Model}
\label{sec:model}

The model is intentionally small, but it is not a claim of expressive
minimality.  Its purpose is to make replay, scheduling, and effects observable
enough to falsify claimed properties.

\subsection{Events, projection, and identity}

An event envelope is a tuple

\begin{equation}
 e = (id, run, seq, kind, cause, emitter).
 \label{eq:event-envelope}
\end{equation}

The event kind is a closed sum containing domain facts, signals, effect
lifecycle events, and evidence facts.  The evidence branch includes
\code{HazardDetected(detail)}; an envelope records it as
\code{EvidenceRecorded(HazardDetected detail)}, and the grader treats that fact
as an explicit blocker.  A log
$L = \langle e_1,\ldots,e_n\rangle$ is ordered and append-only.  The pure
reducer $\mathsf{evolve}:S\times E\rightarrow S$ yields

\begin{equation}
 \project(L) = \mathsf{fold}(\mathsf{evolve}, S_0, L).
 \label{eq:project}
\end{equation}

Event identity is idempotent at the reducer boundary: if an event identifier
has already been applied, a repeated envelope does not mutate state again.  A
duplicate identifier remains an integrity fault for fork assessment, because
idempotent reduction must not silently certify a malformed trace.

State contains domain facts and signals, projected effect descriptors,
evidence facts, integrity faults, and the set of applied event identifiers.
When comparing projected domain state, the implementation excludes the
bookkeeping set of applied identifiers.  This prevents a replay that merely
renames envelopes from being declared semantically different while still
allowing exact-trace equality to detect the change.

\begin{proposition}[Closed-log projection determinism]
For a fixed ordered event log and a pure reducer, repeated evaluation of
Equation~\ref{eq:project} returns the same state.  Oracle calls are not an
exception: their \emph{recorded results}, when present, are events in the fixed
log.
\end{proposition}

This proposition is the trivial base case.  It says nothing about how the log
was produced, whether the event order was unique, whether missing oracle
results can be recovered, or whether effects may be executed again.

\subsection{Behaviors and actions}

A behavior is a tuple

\begin{equation}
 b=(id_b, trigger_b, fire_b, W_b)
\end{equation}

where $trigger_b:S\times E\rightarrow\mathbb{B}$ decides whether an event
enables the behavior, $fire_b:S\times E\rightarrow A^*$ returns actions, and
$W_b$ is a declared write set used only as a diagnostic.  Actions can emit a
domain fact or signal, or request an external effect.  Behaviors cannot emit an
effect outcome or evidence attestation directly; those facts must cross the
interpreter or validation boundary.  This prevents a reducer from minting its
own proof that an external action succeeded or that a run was verified.

The kernel does not execute actions.  Emissions become pending events.  Effect
requests become pending request events and later block until an interpreter
injects a committed, failed, or unknown outcome.  Ordinary .NET tasks, network
clients, mutable caches, and database adapters may exist outside the boundary;
none are reachable from $trigger$ or $fire$.

\subsection{Configurations and scheduling}

A configuration contains projected state, pending events, enabled activations,
outstanding effects, the retained trace, the behavior firing trace, and the next
sequence and generated-identity counters.  A scheduler has two explicit
selectors: one chooses a pending event, and the other chooses an enabled
activation.  Four policies are included: canonical, reverse-event,
reverse-activation, and both reversed.

Writing a configuration as $\langle s,Q,A,O,L,T\rangle$, the two load-bearing
transition rules are shown in Figure~\ref{fig:small-step}.  The helper
$\mathsf{materialize}$ turns a behavior's pure action descriptions into pending
fact/signal events or effect-request events.  It never produces an interpreter
outcome.

\begin{figure}[ht]
\centering
\small
\begin{minipage}{0.96\textwidth}
\[
\frac{
 e=\mathsf{selectEvent}_{\sigma}(Q)\quad
 s'=\mathsf{evolve}(s,e)\quad
 A'=A\mathbin{+}\{(b,e)\mid \mathsf{trigger}_b(s',e)\}
}{
 \langle s,Q,A,O,L,T\rangle
 \rightarrow_{\sigma}
 \langle s',Q\setminus e,A',O,L\mathbin{+}e,T\rangle
}
\quad\textsc{Process-Event}
\]
\[
\frac{
 a=(b,e)=\mathsf{selectActivation}_{\sigma}(A)\quad
 X=\mathsf{fire}_b(s,e)\quad
 (Q',O')=\mathsf{materialize}(Q,O,X)
}{
 \langle s,Q,A,O,L,T\rangle
 \rightarrow_{\sigma}
 \langle s,Q',A\setminus a,O',L,T\mathbin{+}(b,e)\rangle
}
\quad\textsc{Fire-Activation}
\]
\end{minipage}
\caption{Small-step scheduling rules.  Event and activation selection are
explicit.  Emissions become visible to state only after a later
\textsc{Process-Event} step, not within the activation batch for one trigger.}
\label{fig:small-step}
\end{figure}

Processing a pending event first applies it to state and then computes enabled
behaviors from the updated state.  Enabled behaviors fire sequentially, but
their emissions enter the pending queue.  Activations handling the same
triggering event therefore observe the same projected state.  A behavior
activated by a later-processed event observes changes committed by earlier
\textsc{Process-Event} steps, including events materialized from prior
activations.  This cross-event ordering seam mirrors the one found in
ActiveGraph and is sufficient for the read/trigger interference witness in
Section~\ref{sec:confluence}.

Let $C\rightarrow_\sigma C'$ denote one legal event-processing or
activation-firing step under scheduler $\sigma$.  Settlement repeatedly steps
until one of three outcomes is reached:

\begin{itemize}[leftmargin=*]
  \item \textbf{Quiescent}: no pending events, enabled activations, or
  outstanding effects;
  \item \textbf{BlockedAwaitingEffect}: no internal work remains, but at least
  one request has no terminal outcome;
  \item \textbf{Diverged}: a caller-supplied step bound is reached before either
  terminal outcome.
\end{itemize}

The step bound is checked only after testing for a terminal configuration.  A
run that becomes quiescent on the exact final step is therefore not
misclassified as divergence.  Conversely, reaching the bound does not prove a
cycle; the kernel reports \code{StepBoundReached} and makes no stronger claim.
BlockedAwaitingEffect is terminal for the settlement procedure but is not a
quiescent normal form.  The confluence definition in Section~\ref{sec:confluence}
ranges only over executions that reach Quiescent; uniqueness under blocking is
a different property and is not claimed here.

\subsection{Observations and equivalence}

Confluence and fork claims are meaningless without specifying what is observed.
The artifact implements three equivalences:

\begin{enumerate}[leftmargin=*]
  \item \textbf{ExactTrace}: envelope and event lists are structurally equal;
  \item \textbf{ProjectedState}: domain projection is equal after excluding
  applied-event bookkeeping and envelope metadata;
  \item \textbf{FactsOnly}$(K)$: the selected fact keys agree.
\end{enumerate}

The validation additionally reports byte equality when a source artifact
provides it, decision equality when the study records decisions, and structural
prefix equality modulo explicitly excluded lineage fields.  These relations
must not be substituted for one another.  In particular, equal decisions do
not imply equal paths, and equal projected state does not imply an unchanged
external world.

\section{Confluence and Quiescence}
\label{sec:confluence}

The ActiveGraph architecture describes behaviors reacting until quiescence
without a privileged orchestrator.  The working graph is then described as a
deterministic projection of the log.  Projection of a \emph{fixed} log is
deterministic, but the reactive loop can produce different logs under different
legal schedules.  We separate termination from uniqueness of the terminal
projection.

\begin{definition}[Termination at a bound]
For behavior set $B$, initial configuration $C_0$, scheduler $\sigma$, and
bound $m$, $\settle_m(B,C_0,\sigma)$ terminates if it reaches Quiescent or
BlockedAwaitingEffect in at most $m$ steps.  StepBoundReached is evidence of
bounded nontermination, not a proof of infinite reduction.
\end{definition}

\begin{definition}[Schedule confluence under an observation]
A behavior set is confluent for $C_0$ and observation $P$ when every legal
scheduler execution that reaches a quiescent configuration reaches a unique
normal form modulo $\simeq_P$.  Formally, if
$C_0\rightarrow^*_{\sigma_1} C_1$ and
$C_0\rightarrow^*_{\sigma_2} C_2$, and both $C_1$ and $C_2$ are quiescent,
then $C_1\simeq_P C_2$.
\end{definition}

This is Church--Rosser-shaped, but the current artifact does not prove a
Church--Rosser theorem.  It searches scheduler policies, generates restricted
families, and preserves witnesses when the proposition fails.  Newman's lemma
connects termination plus local confluence to confluence in an abstract
rewriting system \cite{newman1942}; our first counterexample removes the
termination premise before the second examines uniqueness of quiescent forms.

\subsection{A1: general termination is false}

\begin{counterexample}[Self-triggering emission]
Let a behavior trigger on signal \code{loop} and emit the same signal.  Seed the
pending queue with \code{loop}.  Every processed signal enables the behavior,
and every firing enqueues another signal.  No quiescent or effect-blocked
configuration is reached.  For every generated bound from 1 through 100, the
artifact reports StepBoundReached at that bound.
\end{counterexample}

The example is intentionally small.  It shows that purity of the reducer and
append-only logging do not imply termination of a reactive runtime.  Detecting
a repeated finite fingerprint would not solve the general problem: generated
identities and sequence counters can make configurations syntactically fresh
while their domain behavior cycles.  The implementation therefore avoids a
heuristic \code{CycleDetected} verdict.

Two natural restoring contracts remain candidates.  One assigns every
internally emitted event a well-founded rank that strictly decreases along
trigger edges.  The other requires an acyclic graph from event kinds to
behavior emissions.  Both would exclude the fixture, but neither is claimed
minimal, and both need refinement for state-dependent triggers.

\subsection{A2: general confluence is false}

\begin{counterexample}[Overlapping writers]
Behaviors $\alpha$ and $\beta$ are both enabled by \code{start}.  Each emits a
write to fact \code{winner}; $\alpha$ writes 1 and $\beta$ writes 2.  Under
canonical activation order the emitted events are processed so that the
terminal winner is 2.  Reversing activation order yields 1.  Both executions
quiesce, but their projected states differ.
\end{counterexample}

FsCheck generalizes this witness over generated integer values by constructing
distinct adjacent values.  Across the generated family, canonical and reverse
activation schedules always disagree.  The write-set diagnostic correctly
reports \code{winner} as a conflict.

It is tempting to restore confluence by requiring declared write sets to be
pairwise disjoint.  That condition works for a deliberately restricted family:
single-trigger behaviors that each unconditionally write a distinct key and do
not read other behaviors' regions or emit cross-triggering signals.  The suite
generates up to eight such writers and checks all four scheduler policies; their
projected states agree.  The structural reason is stronger than those four
samples: every generated key has exactly one unconditional writer, and no
writer reads or triggers another, so any permutation applies the same set of
commuting updates to distinct map entries.  The four policies are executable
regression coverage for that restricted family, not an enumeration of arbitrary
runtime interleavings.

Declared disjointness is not sufficient in general.

\begin{counterexample}[Read/trigger interference with disjoint writes]
Two behaviors, \code{left} and \code{right}, trigger on \code{start} and emit
signals \code{left-ready} and \code{right-ready}.  Their declared write sets are
empty.  A third behavior triggers on \code{left-ready}, reads whether
\code{right-ready} is already in state, and writes the only fact
\code{observed-right}.  All declared write sets are disjoint.  Under canonical
event order the observer records 0; under reverse event order it records 1.
\end{counterexample}

The second witness matters because a write-write conflict detector would
certify it.  The actual dependency crosses four surfaces: an emission, an event
selection, trigger eligibility, and a state read.  Any useful general restoring
condition must constrain at least behavior reads, behavior writes, emitted
event kinds, and the trigger relation.  Under-declared dependencies must either
be rejected or made observable by instrumentation.

\begin{figure}[ht]
\centering
\small
\begin{tabularx}{0.96\textwidth}{@{}p{0.18\textwidth}Y Y@{}}
\toprule
Surface & Canonical event order & Reversed event order \\
\midrule
Emission & \code{left-ready}, then \code{right-ready} & \code{right-ready}, then \code{left-ready} \\
Trigger & Observer becomes eligible before right is projected & Observer becomes eligible after right is projected \\
Read & \code{right-ready} absent & \code{right-ready} present \\
Write & \code{observed-right}=0 & \code{observed-right}=1 \\
\bottomrule
\end{tabularx}
\caption{Counterexample 3 crosses emission, event selection, trigger
eligibility, and a state read despite disjoint declared write sets.}
\label{fig:read-trigger-interference}
\end{figure}

\subsection{Production ordering masks the question}

The ActiveGraph v1.10.0 scheduler uses a FIFO queue, evaluates matching
behaviors in registration order, and dispatches them sequentially.  Work due at
the same delayed tick remains FIFO.  This supplies a reproducible
\emph{ProductionOrder}, and it may be entirely appropriate as an operational
contract.  It is not evidence that alternate legal orders have the same normal
form.

The distinction is practical.  ActiveGraph computes the match set before the
handler loop, so an earlier handler cannot change which behaviors are already
enabled by that trigger.  It does, however, rebuild each handler's view against
the current graph.  An earlier registered handler can emit an event whose
projection is applied synchronously and changes state read by a later handler
of the same original event; behavior matching for the emitted event occurs at a
later FIFO step.  Relation enumeration adds another ordering seam when a
backend query does not specify a canonical order.

The executable kernel deliberately snapshots state across activations enabled
by one trigger, so it does not reproduce this per-invocation production read
timing.  Its tests vary event and activation order directly, and its
read/trigger witness uses distinct later events.  The production seam is a
source-audited limit on generalization, not a property claimed as validated by
the kernel.  A stable FIFO queue can hide either form of dependence in ordinary
replay tests.

\paragraph{Result.}
The strongest supported conclusion is negative and conditional.  General
termination and general confluence are false.  Isolated disjoint writers
converge for the checked family, but declared writes alone do not restore
confluence.  Deriving and validating a general non-interference contract is
future work; this draft does not present a minimal side condition.

\section{Property-Relative Fork Safety}
\label{sec:fork}

A trace fork at cut $k$ retains $L_{<k}$ and gives subsequent events a child
run identity.  For a list representation,

\begin{equation}
  \fork(L,k) = \prefix(L,k).
\end{equation}

That definition makes prefix equality cheap, but it does not by itself make
re-execution or external-world claims sound.  We first describe effects, then
define the properties a fork may be asked to preserve.

\subsection{Effects along three dimensions}

A flat partition into Pure, Idempotent, Cached, and OneShot combines questions
that must be answered separately.  A cached model response describes where
replay material comes from; it does not describe whether the original call had
external cost.  We therefore use three orthogonal dimensions
(Table~\ref{tab:dimensions}).

\begin{table}[ht]
\centering
\small
\begin{tabularx}{\textwidth}{@{}>{\raggedright\arraybackslash}p{0.18\textwidth}>{\raggedright\arraybackslash}p{0.27\textwidth}Y@{}}
\toprule
Dimension & Values & Question \\
\midrule
External footprint & Pure, Idempotent, Compensatable, OneShot, Unknown & What did execution do to the external world, and can that footprint be reconciled? \\
Replay source & Deterministic, Recorded, Uncaptured & Can replay reproduce or serve the result without consulting the live oracle? \\
Lifecycle & Requested, Committed, Failed, Unknown & Is the request/outcome boundary complete and trustworthy at this cut? \\
\bottomrule
\end{tabularx}
\caption{Orthogonal effect dimensions.}
\label{tab:dimensions}
\end{table}

Pure effects are internal descriptions with no external interaction.
Idempotent effects can be repeated or overwritten under an explicit key.
Compensatable effects require a separately verified compensation.  OneShot
effects are irreversible or conservatively treated as non-repeatable.  Unknown
fails closed.  In the validation adapter, model calls are conservatively
OneShot with respect to cost and oracle interaction.  Their outputs may
simultaneously be Recorded, allowing result replay without another call.

Lifecycle is projected from events.  A request must precede exactly one
terminal outcome.  Duplicate requests, conflicting descriptors, an outcome
without a request, or a second outcome after terminal state are integrity
faults.  The first valid request and first valid terminal outcome remain the
projected authority; malformed later events do not overwrite them.

\subsection{Four fork properties}

\begin{definition}[ProjectionReplay]
The child reconstructs the same domain projection as an independently folded
copy of the retained, well-formed, ordered prefix.  No execution occurs.
\end{definition}

\begin{definition}[StrictExecutionReplay]
Every retained effect result needed by execution is deterministic or recorded,
and the cut does not retain a Requested or Unknown lifecycle.  Replaying the
prefix causes zero live oracle or tool calls.
\end{definition}

\begin{definition}[ExternalContinuation]
The runtime or caller supplies a target environment asserted to represent the
retained prefix.  The child may continue after strict replay, provided inherited
committed OneShot effects are served from the record and never executed again,
and no inherited footprint is unknown.  A snapshot identifier, fresh-environment
receipt, or equivalent attestation can discharge the environment premise; when
none is present, the executable verdict is conditional on that premise.  This
property deliberately does not assert that a retrospectively discarded source
suffix left the original world unchanged.
\end{definition}

\begin{definition}[CounterfactualWorld$(\Omega)$]
Let $\Omega$ be a declared set of external observables, such as delivered
messages, provider charges, published records, or payment-ledger entries.  In
addition to ExternalContinuation, the actual world and a world in which the
discarded suffix did not occur agree under every observation in $\Omega$.
Discarded committed OneShot or unknown effects block the claim; discarded
idempotent or compensatable effects make it conditional on reconciliation or
successful compensation.
\end{definition}

The kernel does not inspect the external world.  It checks whether retained
effect evidence is sufficient to license CounterfactualWorld$(\Omega)$ under
the stated footprint-coverage assumption.  Even a Sound verdict is therefore
a contract-relative license, not an empirical observation that two worlds are
identical.

These definitions make ``sound'' property-indexed.  The same cut may be Sound
for ProjectionReplay, Conditional for ExternalContinuation, and Unsound for
CounterfactualWorld.  Table~\ref{tab:forkconditions} summarizes the executable
assessment.

\begin{table}[ht]
\centering
\small
\begin{tabularx}{\textwidth}{@{}p{0.22\textwidth}Y Y@{}}
\toprule
Property & Retained-prefix conditions & Discarded-suffix conditions \\
\midrule
ProjectionReplay & Whole source log is structurally well formed; ordering authority is fixed & None for projected prefix equality \\
StrictExecutionReplay & No open/unknown lifecycle; every needed result is deterministic or recorded & None \\
ExternalContinuation & Strict replay; known footprint; retained committed OneShot effects are not re-executed & Target environment is attested or supplied as a premise \\
\shortstack[l]{CounterfactualWorld\\$(\Omega)$} & ExternalContinuation conditions; effect classification covers $\Omega$ & No unreconciled observed footprint; OneShot or unknown blocks \\
\bottomrule
\end{tabularx}
\caption{Property-relative fork conditions.}
\label{tab:forkconditions}
\end{table}

\subsection{Checked prefix identities}

The executable suite checks the following trace identities over generated logs.

\begin{proposition}[Identity fork]
$\fork(L,|L|)$ retains the full trace.  Under ExactTrace, the shared prefix is
equal to $L$; child lineage is stored outside that shared list.
\end{proposition}

\begin{proposition}[Nested prefix collapse]
For $j\leq i\leq |L|$, retaining $i$ events and then $j$ events has the same
shared prefix as retaining $j$ events directly.  The nested branch still records
the outer child as its parent, so raw branch records are not identical.
\end{proposition}

\begin{proposition}[Projection commutation]
For a well-formed ordered log,
\begin{equation}
 \project(\fork(L,k)) = \project(\prefix(L,k)).
\end{equation}
This is a trace/projection statement, not an external-world theorem.
\end{proposition}

\begin{proposition}[Suffix irrelevance is property-relative]
The discarded suffix is irrelevant to ProjectionReplay of the retained prefix.
It is not generally irrelevant to CounterfactualWorld.
\end{proposition}

The last proposition corrects an over-broad formulation of cheap forking.  A
committed OneShot event in the discarded suffix is a concrete witness: the
forked log excludes it and projects correctly, but the action remains true in
the world.  Likewise, a discarded request without a trustworthy terminal
outcome is Unsound because the external system may have executed it even though
the log never learned the result.

\subsection{Cuts through effects}

For each log, the validator constructs every prefix state in one forward scan
and every environmental suffix finding in one reverse scan.  It then assesses
all $|L|+1$ cuts.  Three boundary cases are especially important.

\begin{figure}[ht]
\centering
\small
\begin{tabularx}{0.96\textwidth}{@{}P{0.20\textwidth}P{0.28\textwidth}Y@{}}
\toprule
Cut & Retained prefix & Licensed claim \\
\midrule
$k_0$ before request & No inherited request & Projection is sound; a later source-side effect may still block a retrospective world claim \\
$k_1$ after request, before outcome & Requested lifecycle only & Strict replay and continuation are refused \\
$k_2$ after recorded commit & Request plus terminal recorded result & Strict replay can be sound; continuation requires zero re-execution \\
\bottomrule
\end{tabularx}
\caption{A single request/outcome pair creates three semantically different
fork regions.}
\label{fig:cut-through-request}
\end{figure}

\paragraph{Before a request.}
The retained prefix does not inherit the effect.  Projection replay is sound.
If the source later committed a OneShot effect, CounterfactualWorld is unsound
for a retrospective fork because that discarded action already occurred.

\paragraph{Between request and outcome.}
The prefix contains a Requested lifecycle.  StrictExecutionReplay and
ExternalContinuation are unsound: the effect may be running, may have committed
without acknowledgement, or may have failed.  A step bound or missing response
must not be treated as a safe cancellation.

\paragraph{After a committed OneShot.}
The response may be Recorded, making strict replay possible with zero external
calls.  ExternalContinuation is Conditional: the interpreter must serve the
recorded result and enforce zero re-executions.  This is a semantic obligation,
not merely an optimization.

These cases expose why a Decider that returns only domain events cannot express
the full agent-runtime question.  The log needs request/outcome boundaries,
footprint evidence, and replay-source evidence before it can license execution
or world-level claims.

\section{Replayability Grades as Evidence Licenses}
\label{sec:grades}

The activegraph-bridge project introduced five replayability levels from
observed through native.  The names arose from implementation failure modes.
We make their meaning explicit by assigning each grade a set of licensed
questions (Table~\ref{tab:grades}).

\begin{table}[ht]
\centering
\small
\begin{tabularx}{\textwidth}{@{}>{\raggedright\arraybackslash}p{0.16\textwidth}>{\raggedright\arraybackslash}p{0.32\textwidth}Y@{}}
\toprule
Grade & Required evidence & Licensed use \\
\midrule
Observed & Some inspectable trace or result & Inspection and available lineage only \\
Envelope & Faithful invocation envelope and completion, no lossy-envelope fact & Recorded-output playback; fork before invocation \\
Boundary & Envelope; mediated effects; clean reconstruction; complete known effect boundaries; no blockers & Re-execution using recorded effects; forks at effect boundaries \\
Checkpointed & Boundary plus recorded checkpoint & Resume from a restorable snapshot \\
Native & Checkpointed plus native-runtime evidence & Native replay, fork, and structural diff \\
\bottomrule
\end{tabularx}
\caption{Replay grades and their licenses.}
\label{tab:grades}
\end{table}

Let $Q(g)$ be the set of questions licensed by grade $g$.  The current taxonomy
is defined as the chain

\begin{equation}
 Q(Observed) \subset Q(Envelope) \subset Q(Boundary)
 \subset Q(Checkpointed) \subset Q(Native).
\end{equation}

This chain is a definitional summary, not a substantive theorem.  The artifact
checks its ordering mechanics, but the paper's claim is narrower: a grade
summarizes the questions supported by retained evidence, while a cut-level
assessment decides one proposed fork.  A richer taxonomy could contain
incomparable evidence sets.

\subsection{Computation from the log}

The grade is computed from projected evidence and effect completeness.  It is
not accepted as a source label.  In particular, a \code{NativeRuntime} marker
alone does not grant Native: the log must also satisfy checkpoint, boundary,
and envelope prerequisites.  Unknown footprint, uncaptured results, incomplete
or unknown lifecycles, integrity faults, lossy envelopes, unmediated effects,
and explicit hazards are blockers.

The adapter follows the same fail-closed rule.  It grants evidence only for
exact recognized source event types.  An event whose name merely contains
\code{verification} remains an unclassified domain signal.  A verification
event must contain an explicit Boolean verdict.  The bridge-specific Boundary
derivation additionally requires an effects-served count and an explicit null
divergence: the source payload's \code{divergence} field must be present and its
JSON value must be \code{null}, asserting that the bridge verifier detected no
unresolved reconstruction divergence.  This is logged evidence, not an
independent proof of equivalence.  These requirements prevent an event-name
coincidence from raising the grade.

\subsection{Grades are not monotone in time}

It is natural to view grades as achievements that only rise.  That view is
incorrect for evidence predicates.

\begin{counterexample}[Hazard downgrade]
A prefix records a faithful envelope, completion, mediated boundary, clean
reconstruction, and successful verification.  It grades Boundary.  Appending a
\code{EvidenceRecorded(HazardDetected detail)} event adds an explicit blocker,
so the extended log grades Envelope.
\end{counterexample}

The grade of a \emph{fixed log or prefix} licenses questions about that evidence.
It need not be monotone under extension.  This matters for forks: the grade of
the shared prefix can be higher or lower than the grade of the completed source
run, and a later hazard must not retroactively be ignored.

\subsection{Connection to fork safety}

The grade is a compact mechanism for checking some fork preconditions, but it
does not replace property-relative assessment.  Boundary evidence suggests
that effect results and reconstruction are available, licensing forks at known
effect boundaries.  A specific cut can still split a request from its outcome.
Likewise, a Boundary run can contain a committed OneShot effect that makes a
CounterfactualWorld claim conditional or unsound.

The validation therefore reports both the run-level grade and all cut-level
verdicts.  A grade answers ``what kinds of questions may this evidence support
in principle?''  A fork assessment answers ``does this particular cut satisfy
the conditions for property $P$?''

\section{Executable Specification and Method}
\label{sec:method}

The artifact is a single .NET solution written in F\#.  Purity is enforced at
the transition boundary, not imposed on the surrounding program.  The source
contains ordinary file I/O, JSON parsing, hashing, mutable library internals,
and command-line orchestration outside the reducer.  There is no database,
dependency-injection container, actor framework, plugin system, network client,
or asynchronous interpreter inside the kernel.

\subsection{Artifact structure}

The artifact includes:

\begin{itemize}[leftmargin=*]
  \item a closed event and effect model;
  \item a pure projector and explicit settlement scheduler;
  \item branch construction with parent/child identity, cut, shared prefix,
  continuation, inherited effects, and prefix hash;
  \item cut-level fork assessments and replay-grade computation;
  \item an ActiveGraph/bridge JSONL normalizer and batch validator;
  \item 50 passing tests, including FsCheck properties and named fixtures;
  \item seven language-neutral JSON conformance vectors and a JSON Schema;
  \item read-only SQLite export and actual-fork audit scripts;
  \item a findings log that records failed checked properties as they are
  discovered.
\end{itemize}

All F\# projects treat warnings as errors.  Pattern matches over closed domain
unions are exhaustive.  External effect execution is not linked into the
kernel, so replay cannot accidentally call a model or tool through an injected
client.

\subsection{Property-based checks}

FsCheck generates bounded event values, writer families, effect dimensions,
and step bounds.  Scheduler policies are enumerated rather than sampled because
the executable model exposes four intended extremes.  Properties include:

\begin{itemize}[leftmargin=*]
  \item deterministic projection and event-identity idempotence;
  \item exact-bound quiescence and bounded self-cycle divergence;
  \item convergence of isolated disjoint writers under all schedulers;
  \item schedule dependence for generated conflicting writers;
  \item identity, nested-collapse, projection, and suffix identities for forks;
  \item verdict matrices across generated footprint and replay-source values;
  \item grade-chain mechanics and prerequisite enforcement.
\end{itemize}

The 50-test suite comprises 15 live FsCheck properties and 35 xUnit facts.
Per-property \code{MaxTest} values range from 100 to 250, totaling 3,300
top-level generated trials in a default full run; properties that enumerate
schedulers or dimension values perform additional checks inside each trial.
No fixed global replay seed is configured.  Regression reproducibility comes
from the checked-in named fixtures rather than from repeating one pseudorandom
stream.

FsCheck shrinking is useful for scalar values, but the semantically important
counterexamples are named and checked in as JSON-backed fixtures.  The fixtures
include overlapping writers, read/trigger interference despite disjoint writes,
a self-cycle, a discarded OneShot effect, an unresolved request, malformed
effect transitions, and forged look-alike evidence types.

Table~\ref{tab:vector-coverage} states the deliberately bounded role of the
seven language-neutral vectors.  They are portability smoke tests, not an
exhaustive cross-product of property, footprint, replay source, lifecycle, and
obligation.

\begin{table}[ht]
\centering
\scriptsize
\begin{tabularx}{\textwidth}{@{}P{0.38\textwidth}Y@{}}
\toprule
Vector & Boundary exercised \\
\midrule
\path{projection-last-write-wins} & ProjectionReplay \code{Sound}; deterministic fact projection \\
\path{discarded-one-shot-is-not-counterfactual} & CounterfactualWorld \code{Unsound}; discarded committed OneShot \\
\path{recorded-effect-strict-replay} & StrictExecutionReplay \code{Sound}; committed Recorded result \\
\shortstack[l]{\code{recorded-unknown-footprint-}\\\code{conditions-external-continuation}} & ExternalContinuation \code{Conditional}; unknown inherited footprint \\
\path{outcome-without-request-is-malformed} & ProjectionReplay \code{Unsound}; lifecycle-integrity blocker \\
\path{boundary-grade} & Boundary grade prerequisites \\
\path{native-grade} & Native grade and prerequisite chain \\
\bottomrule
\end{tabularx}
\caption{Coverage intent of the language-neutral v1 conformance vectors.}
\label{tab:vector-coverage}
\end{table}

\subsection{Real-log normalization}

Validation accepts ordered JSONL.  For SQLite-backed ActiveGraph stores, an
export script uses \code{events.seq} as the ordering authority and emits one
file per run.  Source domain events that are not part of the closed effect or
evidence vocabulary remain signals and are reported by type.  The adapter does
not infer effects from every \code{*.requested} event because real domains use
names such as \code{round.requested}; only a closed list of known external
request types is classified.

This choice creates an explicit adapter-coverage premise.  Unknown footprint
fails closed once an event is classified as an effect, but an unclassified
source event is not automatically treated as an effect.  A Sound continuation
verdict therefore means ``Sound under the profile declaration that the listed
unclassified types are domain-only.''  A new source type that hides an external
interaction can invalidate that verdict.  The validator reports every such type
so the declaration can be audited rather than silently assumed.

For an external request, footprint is read from explicit side-effect metadata.
A write without an idempotency key is OneShot; an unknown tool footprint remains
Unknown.  Model and embedding requests default conservatively to OneShot.  A
response is Recorded only when the log contains an explicit response hash,
output hash, or raw output from which the validator derives a content SHA-256.
The derived hash certifies captured replay material; it is not a provider
receipt or proof of external commitment.

Every source event is assigned a contiguous normalized sequence.  Bridge
events can produce multiple evidence facts; this is why the bridge normalized
count exceeds its source count.  The source event identity is retained for the
first fact, and deterministic derived identifiers are used for additional
facts.  Duplicate or non-contiguous normalized identity remains a blocker.  The
normalizer also retains the cumulative cut after each atomic source event.
These source-boundary positions are reported separately from diagnostic cuts
between co-derived facts.

\subsection{All-cut enumeration}

For a normalized log of length $n$, the validator evaluates cuts
$0,\ldots,n$.  Prefix states are computed with one scan; environmental findings
for suffixes are accumulated in reverse.  It reports both (i) every normalized
position, useful for checking the executable transformation, and (ii) only
source-event boundaries, which are the physically meaningful fork candidates
in the source trace.  The validator reports Sound, Conditional, and Unsound
totals for ProjectionReplay, ExternalContinuation, and CounterfactualWorld,
plus the number of runs containing an unsound counterfactual cut.

Actual ActiveGraph forks require a second audit because the child trace has
lineage and store-level sequence differences.  The SQL audit resolves the
parent cut event, orders the parent prefix and child trace independently, and
compares corresponding event ID, type, actor, payload, frame, cause, and
timestamp.  In this SQLite schema, \code{events.seq} is the database-wide
autoincrement primary key, not a run-local envelope field.  It is used to order
each trace and excluded from equality because copied child events receive new
store positions.  The normalizer then assigns the contiguous per-run sequence
used by Equation~\ref{eq:event-envelope}; that derived ordinal is not a second physical field in the
fork audit.  Child run ID is also excluded.  The audit separately counts
retained external requests and unresolved boundaries using the same explicit
request/outcome vocabulary as the normalizer; unknown event types remain
unclassified rather than being inferred as effects.

\subsection{Reproducibility discipline}

Evidence entries record repository, immutable revision when available, source
path, artifact hashes, counts, parser profile, command, publication status, and
limitations.  Checked-in fixtures are hash-bound and explicitly marked
synthetic or sanitized.  Private and local runtime payloads are not copied into
the repository merely to make a test convenient.  Missing revision or source
provenance is reported as a limitation rather than replaced with a hash-only
claim of reproducibility.

\section{Validation Against Runtime Artifacts}
\label{sec:validation}

We evaluate the contracts against one public trace corpus and one public
post-oracle conformance fork, and retain one local instrumented study as
illustrative evidence.  The public, sealed Synthetic Players capsule contains
an ActiveGraph SQLite store.  The public activegraph-bridge fixture executes an
actual bridge fork after a committed recorded fixture-oracle effect.  The local
executive study records explicit boundary verification but has no public source
remote.  All three contain synthetic experimental content; ``real'' here means
actual runtime artifacts rather than invented shape fixtures.  None is
independent of the ActiveGraph family, and we do not claim human behavioral
validity or real-model performance.

\subsection{Synthetic Players capsule}

The public source is pinned to revision \code{82556ef9dae3}.  The compressed
engine store begins SHA-256 \code{7b7fdfe64b5b}, and the capsule checksum
ledger begins SHA-256 \code{6d80f3b19b92}.  The evidence ledger records the
complete revision and digests.
Running the public verifier reports:

\begin{quote}
CAPSULE VERIFICATION PASS --- 4,919 archived Phase 3--5 runs verified
(4,916 confirmatory + 3 legacy diagnostics).
\end{quote}

The decompressed database contains 5,540 runs and 311,756 events.  Event types
include 36,251 \code{llm.requested}, 36,227 \code{llm.responded}, and 30,397
\code{decision.parsed} records.  The difference between request and response
counts is reflected rather than repaired by the adapter.

The two run counts have different denominators.  The store contains 5,505
completed and 35 incomplete runs.  The public verifier covers 4,916
confirmatory Phase 3--5 runs plus three completed legacy diagnostics.  The
remaining 586 completed runs are outside that replay contract; the 35
incomplete runs are also outside it.  Thus $4{,}919+586+35=5{,}540$.

\subsubsection{Replay grades}

All 5,540 logs grade Observed, and zero are verified \emph{from the run log
alone}.  This does not contradict the passing capsule verifier.  The capsule
attestation is external to each run's event stream.  A grading function that
claims to depend on the log alone cannot silently import that attestation.
This is a useful negative result: sealed-artifact reproducibility and
self-describing run replayability are distinct properties.

\subsubsection{All-cut results}

Table~\ref{tab:syntheticcuts} reports every cut.  The projection total is
$311{,}756 + 5{,}540 = 317{,}296$, because each run of length $n$ has $n+1$
cuts.

\begin{table}[ht]
\centering
\small
\begin{tabular}{@{}lrrr@{}}
\toprule
Property & Sound & Conditional & Unsound \\
\midrule
ProjectionReplay & 317,296 & 0 & 0 \\
ExternalContinuation & 160,076 & 120,969 & 36,251 \\
CounterfactualWorld & 150,230 & 24,172 & 142,894 \\
\bottomrule
\end{tabular}
\caption{All-cut verdicts for 5,540 Synthetic Players logs.}
\label{tab:syntheticcuts}
\end{table}

Every projection cut is sound under the specified domain-projection
observation.  The 36,251 ExternalContinuation-unsound cuts equal the number of
open/cut-through LLM requests.  Direct sequence audit explains the equality:
36,227 requests are immediately followed by their linked response, so each
creates exactly one cut between request and response; the remaining 24 requests
are the final event of failed runs, so each also creates exactly one open cut.
Conditional continuation cuts retain a committed OneShot oracle interaction
and require serving the recorded result without another call.
CounterfactualWorld is stricter: 4,923 runs have at least one Unsound cut
because a discarded oracle interaction already occurred or a boundary is
ambiguous.

The 160,076 Sound ExternalContinuation cuts validate only the
no-inherited-effect case.  A direct per-run ordinal query confirms that every
one lies before the first classified external request, including all cuts in
the 617 runs with no classified request.  The Synthetic Players corpus alone
therefore does not empirically validate a post-oracle continuation,
target-environment attestation, or zero-reexecution receipt.  Generated
properties and named fixtures exercise post-oracle verdicts and typed
no-reexecution obligations; the public conformance fork below separately
exercises their runtime discharge under a fixture trust root.

The validator leaves domain and infrastructure types unclassified, including
\code{behavior.started} and \code{decision.parsed}.  It likewise leaves
\code{infra.*}, \code{round.requested}, and \code{run.completed} unclassified.  In particular,
\code{round.requested} is not misclassified as an external request merely
because its name ends in \code{requested}.  Continuation verdicts are therefore
relative to the audited profile declaration that these reported types are
domain or infrastructure signals rather than hidden external effects.

\subsubsection{Observed forks}

The SQLite run table contains 121 child runs from 41 distinct parents and no
nested child in this sample.  All forks cut at a \code{round.played} event.  The
audit compares 16,625 parent-prefix events with the beginning of their child
traces and finds zero mismatches in event ID, type, actor, payload, frame, cause,
or timestamp.  Every cut event resolves.  Run identity and \code{events.seq}
are excluded because they are expected to differ: in this store schema,
\code{seq} is the database-wide autoincrement key used for ordering, not a
run-local field.  The per-run normalized sequence is derived after export and
was therefore not an additional physical field available for comparison.

None of the 121 retained prefixes contains a classified external request, and
none contains an unresolved request.  The observed forks therefore satisfy the
current ExternalContinuation precondition for a domain-only cut family.  This
is strong evidence for cheap structural branching where the runtime actually
forked.  It is not evidence for an effect-boundary fork: the observed sample
does not contain one.

\subsection{Public post-oracle conformance fork}

The companion activegraph-bridge artifact is public at revision
\code{8855d3a9e779}.  The complete
\href{https://github.com/yoheinakajima/activegraph-bridge/tree/8855d3a9e779362f713b08bceb58d7d5db671c7d/evidence/post-oracle-fork-v1}{public post-oracle fixture and verifier}
are revision-pinned.
It records an actual bridge child fork after a committed effect classified as
OneShot footprint plus Recorded replay source.  The source run makes one
deterministic offline fixture-oracle call.  Verification and the child serve
the committed outcome from the record, and the child diverges only after a
changed tool result.

The receipt binds the parent and child run identities, cut boundary, parent log
hash, exact retained-prefix hash, child log hash before receipt emission,
inherited request and outcome identities, response hash, source and target
runtime fingerprints, and target-environment claims.  Its verifier checks those
hashes, the zero-call counter, and an HMAC-SHA256 signature under a configured
trust root whose claims include the child, cut, prefix, and target fingerprint.
The published command reports:

\begin{quote}
POST-ORACLE FORK RECEIPT PASS --- committed oracle served from record;
0 inherited external calls.
\end{quote}

The parent JSONL digest begins \code{884e8ae34296}; the child begins
\code{747305a6774a}.  The receipt file begins \code{da5503a80b92}, and its
canonical internal receipt hash begins \code{5d36df7dc177}.  The verified receipt names inherited request
\code{evt\_008}, its recorded outcome \code{evt\_009}, one source oracle call,
and zero inherited external calls in the fork.

This discharges the prior implementation gap for a public conformance case.  An
inherited recorded oracle outcome can cross a fork boundary without a second
call, and the target environment premise can be authenticated against a
configured trust root.  The oracle is deterministic and offline, and the HMAC
key is deliberately published as a fixture key.  The result does not attest a
real provider, production identity, model quality, or production environmental
fidelity.

\subsection{Illustrative local executive study}

The second artifact contains 12 synthetic vendor due-diligence cases evaluated
under two deterministic offline mock strategies, for 24 runs and 2,592 source
events.  The manifest states that no real LLM or external provider was used and
limits the claim to instrumentation feasibility.  It records fresh-factory
reconstruction and enables boundary verification.

The artifact files are hash-bound.  Hash prefixes for the run store, manifest,
metrics, and paired results are \code{7a38a1e8}, \code{da1cbc3b},
\code{9f9242e1}, and \code{f4d95dbb}, respectively; the evidence ledger records
the complete hashes.  Review exposed that the containing migration-lab had no
commit or remote.  We repaired the local half of that defect: source revision
\code{54dfde4d7379} and release-bundle revision \code{201efa7} bind a redacted
manifest, metrics, pairs, report, compressed run store, and verifier.  The
repository still has no configured remote.  This result remains illustrative,
is excluded from the abstract, and is not externally source-reproducible until
those revisions are published.

\subsubsection{Decision and path equivalence}

The paired artifact reports 66.6667\% final-decision agreement and 0\% ordered
tool-path agreement.  Eight of twelve pairs reach the same decision via
different paths, yielding 66.6667\% same-decision/different-path.  This is not a
model-quality result---both strategies are deterministic mocks.  It is an
instrumented witness that decision equivalence and trace equivalence answer
different questions.  A system that reports only agreement can hide total path
divergence; a system that requires exact paths can reject decisions that are
equal under a task-level observation.

\subsubsection{Grades and cuts}

Each run's source events normalize to two additional evidence facts: fresh
reconstruction and mediated boundary.  The 2,592 source events therefore become
2,640 normalized events.  All 24 logs grade Boundary and all 24 are verified
from explicit log facts.  The 48 additional facts create 48 diagnostic
intra-source positions that cannot be selected as cuts in the atomic source
log.  We therefore report source-event-boundary verdicts as the primary bridge
result and retain every-normalized-position counts as a transformation check.

\begin{table}[ht]
\centering
\small
\begin{tabular}{@{}lrrr@{}}
\toprule
Property & Sound & Conditional & Unsound \\
\midrule
ProjectionReplay & 2,616 & 0 & 0 \\
ExternalContinuation & 1,968 & 264 & 384 \\
CounterfactualWorld & 0 & 264 & 2,352 \\
\bottomrule
\end{tabular}
\caption{Source-event-boundary verdicts for 24 bridge logs.}
\label{tab:bridgecuts}
\end{table}

All 2,616 source-boundary ProjectionReplay cuts are Sound; the diagnostic over
every normalized position is 2,664 Sound.  For completeness, the normalized
ExternalContinuation totals are 1,992 Sound, 288 Conditional, and 384 Unsound;
the CounterfactualWorld totals are 0 Sound, 288 Conditional, and 2,376 Unsound.
Every run contains a OneShot
\code{submit\_recommendation} effect.  No cut is unconditionally Sound for
CounterfactualWorld: before the commit, the effect lies in the discarded
suffix; after the commit, continuation is Conditional on not executing it
again.  The result demonstrates why Boundary grade licenses effect-boundary
operations but cannot promise a world in which discarded actions never
happened.  These counts use the adapter-wide $\Omega_{profile}$, under which
\code{submit\_recommendation} is in scope.  The public bridge v0.2 receipt
records per-effect observable tags, but the legacy executive-study adapter does
not expose such tags to this analysis.  A narrower observable set could change
the world verdicts; we do not report that unmeasured result.

\subsection{Validation status}

\paragraph{V1: compute grade distributions.}
Completed.  The public capsule yields Observed for 5,540 logs because
verification is external to the log; the bridge yields Boundary for 24 logs
because the necessary evidence is explicit.  The contrast supports log-derived
grading rather than profile-derived labels.

\paragraph{V2: inspect real forks and their preconditions.}
Completed for the public capsule and the public conformance fixture.  All 121
observed corpus forks preserve the compared prefix fields and retain no
classified external request.  The companion fixture tests one actual bridge
fork after a committed recorded oracle call and verifies zero inherited
external calls plus an authenticated fixture-environment premise.  It is
conformance evidence, not a real-provider or production-attestor result.

\paragraph{V3: test production confluence conditions.}
Partially completed.  Source audit confirms that ActiveGraph imposes FIFO plus
registration order, which can mask non-confluence.  The property suite bypasses
that order.  A complete inventory of production behavior reads, writes,
emissions, triggers, and relation-order dependencies was not available, so we
cannot certify or falsify a general production non-interference contract.

\paragraph{Combined interpretation.}
The evidence supports deterministic projection, structurally correct observed
prefix copying, and one public post-oracle continuation under a verified
conformance environment.  It does not support an unrestricted claim that any
historical cut is safely executable or environmentally counterfactual.  Those
stronger claims require effect-complete evidence and a property-specific
precondition.

\section{Related Work and Novelty Boundary}
\label{sec:related}

\paragraph{Reducer architectures.}
Mealy machines produce output and a next state from a current state and input
\cite{mealy1955method}.  Elm's update loop maps a message and model to a new
model and command \cite{elmarchitecture}.  Chassaing's functional event-sourcing
Decider separates decision and evolution and demonstrates composition in F\#
\cite{chassaing2021decider}.  These precedents invalidate a novelty claim for
the reducer shape.  Our contribution begins where the runtime distinguishes
replayable facts from externally executed effect requests and asks about
specific fork cuts.

The title ``Agent Algebra'' was also abandoned because Burch, Passerone, and
Sangiovanni-Vincentelli already use Agent Algebra for a general framework of
concurrent models, refinement, composition, and conformance
\cite{burch2003agent}.  The present project is consequently named for the laws
it checks rather than an algebra it does not originate.

\paragraph{Agent calculi and mechanized semantics.}
$\lambda_A$ gives a typed lambda calculus for agent composition with oracle
calls, bounded fixpoints, probabilistic choice, and mutable environments.  It
proves type safety and bounded termination in 1,519 lines of Coq with 42
theorems \cite{liu2026lambdaa}.  Its properties are primarily about
composition-time well-formedness and bounded constructs.  We do not duplicate
that contribution; we study runtime event logs, schedule-sensitive settlement,
and effect-boundary forks.

LLMbda models agent conversations, fork and clear operations, code generation,
and labeled information flow.  Its central result is a machine-checked
termination-insensitive probabilistic noninterference theorem
\cite{garby2026llmbda}.  Our work does not claim noninterference or prompt-
injection safety.  Its use of conversation forks is related operationally, but
our property of interest is whether an event-log cut licenses replay and an
external-world claim.

Schlapbach formalizes Schema-Guided Dialogue and the Model Context Protocol in
a process calculus and studies structural bisimilarity and expressiveness
\cite{schlapbach2026protocols}.  That is a protocol-level relation.  We treat
model and tool interactions as effect boundaries inside a runtime trace and do
not compare protocol expressiveness.

\paragraph{Runtime substrates and governance.}
Agent libOS supplies a capability-controlled runtime with process identity,
memory, skills, JIT tools, checkpoints, restore, fork, and commit
\cite{zhang2026libos}.  It demonstrates operational mechanisms and a safety
benchmark.  We do not contribute a capability system.  Our narrower question
is which algebraic and evidentiary conditions make a particular event-log fork
sound for a named observation.

McCann's Effect-Transparent Governance formalizes handler-mediated effect
governance in Rocq and proves semantic preservation, expressive minimality, and
decidability boundaries \cite{mccann2026governance}.  This directly motivates
our decision to concede minimality.  We use an effect-transparent boundary as
an executable assumption, then investigate replay-source, lifecycle, and
external-footprint evidence rather than proving governance expressivity.

AgentSpec lets users state runtime rules as triggers, predicates, and
enforcement mechanisms and evaluates safety across code, embodied, and driving
domains \cite{wang2026agentspec}.  Its contribution is customizable runtime
enforcement.  Our laws do not enforce domain policy; they decide which replay
or fork statements are licensed by a retained trace.

\paragraph{ActiveGraph.}
ActiveGraph makes the event log authoritative and presents deterministic
replay, cheap forks, and lineage as architectural properties
\cite{nakajima2026log}.  This paper is explicitly a follow-on.  It preserves
the fixed-log replay result, shows that reactive settlement needs separate
termination and confluence conditions, and reframes cheap forking as a semantic
claim indexed by projection, execution, continuation, or external-world
equivalence.  The public ActiveGraph implementation and its downstream
artifacts supply the validation substrate unavailable to purely formal work.

\paragraph{What remains distinct.}
We are not aware, within the cited agent-calculus and runtime work, of an
executable study that jointly (i) permutes runtime firing order, (ii) preserves
negative confluence witnesses, (iii) partitions effect evidence along
footprint/source/lifecycle dimensions, (iv) assesses every cut under multiple
fork properties, and (v) checks those laws against an existing event-sourced
runtime's archived traces and observed forks.  This is the bounded novelty claim
for the draft; it should be revisited with a broader systematic search before
submission.

\section{Discussion}
\label{sec:discussion}

\subsection{``The log replayed'' is not one assertion}

The empirical results expose a hierarchy of claims that are often collapsed in
agent-runtime descriptions.

\begin{enumerate}[leftmargin=*]
  \item \textbf{Byte replay}: a serialized artifact is identical under a
  verifier.
  \item \textbf{Trace replay}: the same ordered semantic events are retained.
  \item \textbf{Projection replay}: folding the events yields an equivalent
  state under a chosen observation.
  \item \textbf{Path replay}: the same model/tool/action sequence occurs.
  \item \textbf{Decision equivalence}: a task-level conclusion agrees.
  \item \textbf{Execution replay}: the run can be reproduced with zero live
  calls by serving deterministic or recorded effect results.
  \item \textbf{Environmental equivalence}: the external world matches the
  world described by the retained prefix.
\end{enumerate}

None of these relations automatically implies all the others.  The bridge
study has equal decisions in two thirds of pairs and equal ordered paths in
none.  The Synthetic Players capsule passes an external byte-exact verifier but
its individual logs contain no verification attestation, so they remain
Observed under a log-alone grade.  Every prefix projects, but cuts through
requests fail execution continuation.  These are not contradictions; they are
different predicates.

\subsection{Cheap fork is semantic before it is fast}

Sharing an immutable prefix is computationally cheap.  The harder question is
what the child is allowed to do next.  A runtime that advertises \code{fork(k)}
should specify at least:

\begin{itemize}[leftmargin=*]
  \item the observation under which parent prefix and child are equivalent;
  \item the ordering authority for the prefix;
  \item whether a cut may split a request from its outcome;
  \item how recorded results are served without re-execution;
  \item which retained external effects are inherited;
  \item whether the child receives an isolated/snapshotted environment or the
  already-mutated source environment;
  \item the treatment of external effects in a retrospectively discarded
  suffix.
\end{itemize}

Without these clauses, ``cheap'' describes storage cost while leaving semantic
cost undefined.  With them, the implementation can refuse unsafe cuts, require
an idempotency key, request compensation, select a snapshot, or downgrade the
claim to ProjectionReplay.

\subsection{Deterministic ordering can conceal non-confluence}

FIFO and registration order are useful engineering contracts.  They make one
execution repeatable and improve debuggability.  But they can turn an accidental
ordering dependency into a stable behavior, making ordinary replay tests pass
while schedule independence remains false.  This matters when a runtime later
parallelizes handlers, changes relation storage, distributes a queue, or moves
from one database backend to another.

The appropriate response need not be to demand full confluence.  Some systems
may intentionally make order semantic.  The contract should then name the
canonical order, preserve it in the log, and reject backends that cannot supply
it.  Systems that want schedule independence need a stronger behavior
non-interference discipline.  The current work distinguishes the choices but
does not prescribe one universally.

\subsection{The log should carry its own licenses}

The gap between capsule verification and log-alone grade suggests a design
principle: evidence needed to authorize later replay or fork should be retained
with the run or its signed sidecar.  A minimum receipt should include canonical
sequence, run identity, parent and cut lineage, effect request identity,
footprint annotation, outcome linkage, response content hash, commitment or
failure evidence, reconstruction method, and verification result.

This does not require every event to be trusted merely because it is present.
The event type must be recognized, the producer or signature may need
authentication, and cross-event integrity must be checked.  The present
artifact enforces structural integrity and exact recognized types but does not
implement cryptographic attestation.  A future signed evidence capsule could
make grade computation portable across implementations.

\subsection{The graph may be a projection, but the scheduler remains real}

The original minimality experiment asked whether an agent graph is fundamental
or merely a projection over event-producing reducers.  The present findings
support part of that intuition: replayed domain state can indeed be treated as
a projection.  They also reveal what the reduction omits.  Once multiple
behaviors react, event and activation order become part of the operational
semantics unless a confluence condition removes them.  Once effects leave the
process, the external world becomes a second state that the event projection
can describe but not reverse.

Thus reducing the runtime to a pure fold is useful but incomplete.  The deeper
object is an event-producing transition system plus a scheduler, an effect
interpreter, and evidence linking interpreter outcomes back to the log.  The
graph may be derived; the ordering and effect contracts cannot be wished away.

\subsection{Operational recommendations}

The results suggest concrete contracts for event-sourced agent runtimes:

\begin{enumerate}[leftmargin=*]
  \item Treat production order as an explicit scheduler policy and test at
  least event-order and activation-order perturbations.
  \item Record behavior read, write, emission, and trigger dependencies if
  schedule independence is a goal; a write set alone is insufficient.
  \item Separate replay source from external footprint.  ``Cached'' does not
  mean ``never happened'' or ``safe to repeat.''
  \item Make request and terminal outcome a one-to-one integrity boundary.
  Refuse a cut through Requested or Unknown without an owner-selected recovery
  policy.
  \item Give every fork API a property parameter or an equivalently explicit
  mode.  Do not return a single unqualified \code{safe} Boolean.
  \item Preserve verification and reconstruction evidence inside a portable,
  hash-bound run capsule.
  \item Distinguish actual observed forks from every-cut counterfactual tests.
  The former show where users forked; the latter reveal untested hazards.
\end{enumerate}

These recommendations are compatible with F\#, C\#, Python, Rust, or
TypeScript.  The conformance vectors intentionally encode events and expected
verdicts in language-neutral JSON so another implementation can challenge the
semantics.

\section{Limitations and Threats to Validity}
\label{sec:limitations}

\paragraph{No mechanized proof.}
The propositions and contracts are encoded as executable properties and checked over finite
generators and artifacts.  A passing FsCheck property is not a proof over all
behaviors, schedulers, or logs.  Future work may mechanize selected definitions
after the operational contracts stabilize.

\paragraph{No minimal confluence condition.}
We falsify general confluence and declared-write disjointness.  We check a
restricted isolated-writer family.  We do not establish a necessary or
sufficient general condition over reads, writes, emissions, triggers, and state
domains.  Calling the current restricted condition ``the'' restoring condition
would overstate the result.

\paragraph{Limited scheduler space.}
The executable suite uses four deterministic scheduler extremes.  They expose
the preserved counterexamples but do not enumerate every interleaving for an
arbitrary behavior set.  Relation ordering is discussed from source audit but
not modeled as a separate production activation dimension.  Seeded randomized
and bounded exhaustive interleaving exploration remain future checks.

\paragraph{Adapter assumptions.}
Real-log classification relies on explicit event names and payload fields.  The
adapter is deliberately conservative, but mistakes remain possible.  Model
calls are treated as OneShot for cost/oracle semantics; another property might
classify a read-only deterministic local model differently.  A derived content
hash proves captured bytes, not provider authenticity, commitment, or
idempotence.

The v0 effect descriptor has no per-effect observable-set label.  Its
CounterfactualWorld result therefore uses the selected profile's entire
recognized external surface as $\Omega_{profile}$.  Narrower claims require
adapter-supplied observable tags and new conformance vectors.

Unknown footprint fails closed only after the adapter recognizes an event as an
effect.  Unclassified source events remain reported signals.  The all-cut Sound
counts therefore depend on an audited profile declaration that those event
types do not hide external interactions.  A newly introduced effect type can
make an old verdict false until the adapter is updated.

\paragraph{Partial failure and compensation scope.}
Failed remote operations may have partially committed before reporting an
error.  The assessor conservatively conditions or blocks discarded failures by
footprint, but it cannot inspect the remote system.  Likewise, successful
compensation does not erase a temporal exposure window or third-party secondary
effects.  A caller whose $\Omega$ includes those histories needs stronger
reconciliation evidence than the current effect descriptor records.

\paragraph{Whole-log malformedness policy.}
The current assessor treats structural malformedness anywhere in the source log
as a blocker for every cut, including projection.  This is fail-closed and
appropriate for the evidence study, but it is stricter than a theorem about an
independently authenticated well-formed prefix.  The projection suffix identity is
therefore stated only for well-formed ordered logs.

\paragraph{Synthetic content.}
The Synthetic Players capsule contains archived experimental runs, but the
present analysis concerns runtime traces rather than human ground truth.  The
bridge study uses deterministic offline mocks and explicitly makes no real-model
performance claim.  Its agreement numbers illustrate equivalence relations,
not predictive validity.

\paragraph{Fork coverage.}
The 121 observed forks are useful positive evidence, but all cut at
\code{round.played} and retain no classified external request.  No observed
fork in the public capsule crosses a recorded oracle or irreversible effect
boundary.  The difficult cases are exercised by all-cut analysis and fixtures,
not by an actual public fork with a no-reexecution receipt.  Moreover, all
160,076 Sound public ExternalContinuation cuts precede the first classified
effect.  The current target-environment premise is caller-supplied and has no
implemented attestation-discharge path, so neither observed forks nor those
Sound counts validate one.

\paragraph{Bridge provenance.}
The bridge artifact is locally content-addressed and now bound to local source
and release-bundle commits with a standalone verifier.  The repository still
has no public remote, so a fresh external reviewer cannot fetch those revisions.
The result remains illustrative and excluded from the abstract until the
bundle is published.  This limitation is particularly important because
provenance is part of the paper's thesis.

\paragraph{Production behavior inventory.}
The ActiveGraph source establishes FIFO and registration-order scheduling, but
we did not recover a complete deployed behavior dependency inventory.  We
therefore cannot say whether production behaviors satisfy a future
non-interference condition.  Static source inspection of the scheduler is not
production confluence validation.

\paragraph{No external runtime corpus.}
The public corpus, observed forks, bridge implementation, and local executive
study all descend from ActiveGraph.  The language-neutral vectors support other
implementations, but no LangGraph, Temporal, or other independent runtime trace
has yet been normalized.  Generalization beyond the source family is therefore
conceptual rather than empirically demonstrated.

\paragraph{No authorization or information-flow result.}
Effect descriptors classify replay and environmental conditions; they do not
grant capabilities or prove that a caller is authorized.  The artifact also
does not track secrecy labels or prove noninterference.  Those topics are
covered by distinct systems and calculi and remain explicit non-goals.

\paragraph{No performance benchmark.}
The all-cut implementation is efficient enough for the reported corpora, but
we make no throughput, storage, or latency claim.  ``Cheap'' in this paper is
analyzed semantically.  A separate benchmark would be needed for production
performance.

\section{Conclusion}
\label{sec:conclusion}

Deterministic projection of a fixed log is valuable, but it is the beginning of
an agent-runtime replay contract rather than the end.  Reactive settlement can
diverge, a stable production order can mask non-confluence, and disjoint
declared writes do not exclude schedule-sensitive reads or trigger
interference.  A fork can preserve its trace prefix while failing to reproduce
execution or the external world.

The executable contract makes those distinctions checkable.  Effect
footprint, replay source, and lifecycle answer separate questions.  Fork safety
is indexed by an observation: projection, strict execution replay, continuation
in a prefix-reconstructed environment, or counterfactual-world equivalence.
Replay grades summarize retained evidence, while cut-level assessment decides
whether a particular branch meets the property-specific preconditions.

The runtime artifacts support both positive and negative conclusions.  Across
5,540 public capsule logs, every one of 317,296 projected prefixes is sound
under the stated projection observation.  Across 121 observed forks, 16,625
shared-prefix events show zero structural mismatches modulo expected lineage
fields, and the observed corpus cuts avoid external requests.  Yet 36,251 hypothetical
cuts split model requests and are unsound for continuation, while thousands of
discarded oracle interactions prevent an unchanged-world claim.  A separate
public offline bridge fixture crosses one committed recorded oracle boundary
and directly observes no second fixture-oracle call; its verifier checks
receipt/log consistency and a fork-bound caller assertion under a public
conformance key.  A local executive study further illustrates that decision
and path equivalence can separate, but it remains outside the public headline
evidence until its commit-bound bundle is published.

The practical conclusion is concise: replay is a family of assertions, and a
fork is sound only relative to one of them.  Event-sourced agent runtimes should
record the evidence that licenses each assertion, expose ordering assumptions,
and refuse to turn a cheap prefix copy into an unqualified promise about
execution or the world.


\bibliographystyle{plain}
\bibliography{references}

\end{document}
